% ==============================================================================
% SECCIÓN 3: RECURSOS SOFTWARE DE DESARROLLO
% ==============================================================================
\section{Recursos Software de Desarrollo}

La aplicación se desarrollará en \textbf{Flutter}. Trabajaremos principalmente con el \textbf{IDE Android Studio} con las extensiones de Dart y Flutter, por las facilidades que nos brinda a la hora de hacer tests unitarios y probar los prototipos en los propios dispositivos físicos.

% Revisar las tecnologías utilizadas para el backend
\subsection{Ecosistema y Herramientas de Desarrollo}
\begin{itemize}[leftmargin=*]
  \item \textbf{Framework de Desarrollo Frontend:} \textbf{Flutter SDK (Dart)}. Permite generar aplicaciones nativas multiplataforma desde un único código fuente, asegurando un rendimiento fluido en las tabletas.
  
  \item \textbf{Entorno de Desarrollo (IDE de Flutter):} \textbf{Visual Studio Code / Android Studio} con las extensiones oficiales de Flutter y Dart. Ofrece:
    
    \begin{itemize}
      \item \textit{Hot Reload} y \textit{Hot Restart} para modificar y personalizar componentes gráficos y estilos visuales en tiempo de ejecución de manera instantánea.
      
      \item Suite integrada para ejecución automatizada y depuración de pruebas unitarias (*Flutter Test*).
      
      \item \textit{Flutter DevTools} para análisis de rendimiento, inspección del árbol de widgets y perfiles de memoria.
    \end{itemize}
  
  \item \textbf{Backend y API REST:} \textbf{Python con FastAPI}, aportando tipado seguro, arquitectura desacoplada y alta escalabilidad en la gestión de servicios web.
  
  \item \textbf{Persistencia y Base de Datos:} \textbf{PostgreSQL} para almacenar y gestionar los datos de la aplicación.

  
  \item \textbf{Catálogo de Pictogramas:} Catálogo oficial de pictogramas de \textbf{ARASAAC} (licencia Creative Commons) para comunicación aumentativa.
  
  \item \textbf{Diseño UX/UI:} \textbf{Figma} para la confección y testeo de wireframes accesibles previos a la fase de programación.
\end{itemize}
